\documentclass[twoside]{article}

\usepackage[preprint]{aistats2027}
\usepackage[round,authoryear]{natbib}
\usepackage{amsmath,amssymb,amsthm}
\usepackage{mathtools}
\usepackage{microtype}
\usepackage{booktabs}
\usepackage{url}

\renewcommand{\bibname}{References}
\renewcommand{\bibsection}{\subsubsection*{\bibname}}
\bibliographystyle{apalike}

\newtheorem{theorem}{Theorem}
\newtheorem{lemma}[theorem]{Lemma}
\newtheorem{proposition}[theorem]{Proposition}
\newtheorem{corollary}[theorem]{Corollary}
\newtheorem{definition}[theorem]{Definition}
\newtheorem{remark}[theorem]{Remark}
\newcommand{\R}{\mathbb R}
\newcommand{\E}{\mathbb E}
\newcommand{\cA}{\mathcal A}
\newcommand{\cD}{\mathcal D}
\newcommand{\cL}{\mathcal L}
\newcommand{\cS}{\mathcal S}
\newcommand{\conv}{\operatorname{conv}}
\newcommand{\spanop}{\operatorname{span}}
\newcommand{\rank}{\operatorname{rank}}
\newcommand{\supp}{\operatorname{supp}}

\begin{document}
\runningauthor{}

\twocolumn[
\aistatstitle{One Short of Steinitz: Weighted Regression Selection}

\aistatsauthor{Pahan Dewasurendra}

\aistatsaddress{Johns Hopkins University}
]

\begin{abstract}
How well can least squares perform on a full dataset after training on a
small weighted subset? A recent COLT open problem determined the exact
worst-case ratio for budgets at most the dimension and at least twice the
dimension, leaving every intermediate budget open and reporting an
unpublished endpoint argument. We give a complete proof that, for
$d$-dimensional homogeneous linear regression and the minimum-norm ERM, the
exact ratio with $2d-1$ selected examples is $1+1/d$.

The proof identifies the only obstruction to exact recovery one example
below the Steinitz threshold. Unless the nonzero full-loss gradients lie on
$d$ independent lines and every useful zero-gradient feature is absent,
$2d-1$ examples recover the full ERM exactly. On the exceptional cross
configuration, the problem splits into one-dimensional regressions and
omitting one member of the least-loss line costs at most a $1/d$ fraction
of the optimum. A matching construction handles the minimum-norm tie rule.

We also give a harmonic lower bound for every intermediate budget and show
that it is exact for rank-additive circuit decompositions. This class
contains the extremal construction and isolates the remaining obstacle for
the full open curve: overlapping positive circuits can only be handled by
a new reduction, not by ordinary volume sampling.
\end{abstract}

\section{Introduction}

Data selection asks a deliberately data-centric question. Fix the training
algorithm, reveal the entire finite dataset, and choose a small subset on
which to train. How close can the resulting model come to the model trained
on all examples? This differs from feature selection, random subsampling,
and label-efficient active learning. The selector here knows both features
and responses, may assign nonnegative weights, and is evaluated on the
original observed dataset.

Even homogeneous least squares has a sharp unresolved regime. Let
$\cA_*$ return the minimum Euclidean-norm empirical risk minimizer. For a
dataset in $\R^d\times\R$, \citet{hanneke2025data} proved that the worst
multiplicative loss ratio is infinite below $d$ selected examples, equals
$d+1$ at $d$, and equals one at $2d$ or more. Their COLT open-problem note
asks for the exact value between $d$ and $2d$. The authors report belief in
an unpublished, ad hoc argument for the value $1+1/d$ at $2d-1$
\citep{hanneke2025open}. The two proved endpoints use different geometric
tools. Size-$d$ volume sampling gives $d+1$
\citep{derezinski2017unbiased,derezinski2018reverse}, while Steinitz's
theorem gives exact recovery at $2d$ \citep{steinitz1916bedingt}.

We give a complete endpoint proof. The argument does not interpolate an
existing volume-sampling estimate. Instead, it characterizes when the
$2d$-point Steinitz guarantee cannot be improved by one point. In that
case all informative gradients form a cross, namely $d$ pairs of opposite
directions. The regression objective then separates over those directions.
A one-dimensional averaging argument chooses which pair to leave
incomplete and gives the factor $1+1/d$.

Our contributions are:
\begin{itemize}
\item We prove the exact universal identity
\begin{equation}
  \sup_D \frac{L_D^*(2d-1;\cA_*)}{L_D^*}=1+\frac1d.
  \label{eq:headline}
\end{equation}
This includes rank-deficient datasets, zero optimal loss, arbitrary
dataset size, and the minimum-norm tie rule.
\item We characterize the datasets whose full optimum needs $2d$ weighted
examples to become a \emph{unique} selected ERM. Their nonzero gradients
lie on $d$ independent lines, both orientations occur on each line, and
every zero-gradient feature is zero. We separately state the exact
criterion for the minimum-norm rule, which can exploit a rank-deficient
objective.
\item For every budget $d+k$, we give an explicit harmonic lower bound.
We prove the same expression as an upper bound whenever the normalized
loss admits a rank-additive conformal circuit decomposition. Direct sums
of balanced simplex blocks are extremal in this class.
\end{itemize}

The last result is deliberately not presented as a solution of the full
intermediate curve. Positive-basis structure permits overlapping circuits
\citep{davis1954positive,schoch2020simplicial,hare2021positive}. Our exact
argument covers rank-additive circuits, and computations suggest overlap
helps, but a universal reduction remains open.

\paragraph{Relation to other subset problems.}
Strong regression coresets preserve the objective uniformly over all
parameters, which is a different and generally more demanding target
\citep{boutsidis2013coresets,feldman2020coresets}. Active regression chooses
labels without seeing all responses \citep{chen2019active}. Machine
teaching also optimizes training examples for a fixed learner, but commonly
allows synthesis or targets parameter identification rather than the
full-dataset loss \citep{ma2018teacher}. The present formulation follows
\citet{hanneke2025data}: labels are known, selected examples must come from
the dataset, and the learner is fixed.

\section{Problem and Geometry}

Let $D=\{(x_i,y_i)\}_{i=1}^N$ be a finite multiset in
$\R^d\times\R$. Define
\begin{equation}
 L_D(w)=\frac1N\sum_{i=1}^N(x_i^\top w-y_i)^2,
 \qquad L_D^*=\min_w L_D(w).
\end{equation}
For a convex combination $F=\sum_i\lambda_i\ell_i$, where
$\lambda_i\geq0$, $\sum_i\lambda_i=1$, and
$\ell_i(w)=(x_i^\top w-y_i)^2$, let $\cA_*(F)$ be its minimum-norm
minimizer. The best loss using at most $n$ examples is
\begin{equation}
 L_D^*(n;\cA_*)=
 \inf_{\substack{\lambda\geq0\\|\supp(\lambda)|\leq n}}
 L_D\!\left(\cA_*\!\left(\sum_i\lambda_i\ell_i\right)\right).
 \label{eq:selected}
\end{equation}
Repeated examples and zero weights make this equivalent to the original
definition with $n$ selected examples. Ratios use the conventions of
\citet{hanneke2025data}: $0/0=1$ and a positive numerator over zero is
infinite.

Assume first that the feature matrix has rank $d$. The full minimizer
$w_*$ is unique. Put
\begin{equation}
 r_i=y_i-x_i^\top w_* ,\qquad g_i=r_i x_i,
 \qquad H=\frac1N\sum_i x_ix_i^\top.
 \label{eq:residual-gradient}
\end{equation}
The normal equations and Pythagorean identity are
\begin{equation}
 \sum_i g_i=0,\qquad
 L_D(w_*+v)=L_D^*+v^\top Hv.
 \label{eq:pythagorean}
\end{equation}
The vectors $g_i$ are the individual gradients at $w_*$ up to a common
nonzero scalar and sign.

The tie rule enters through one elementary criterion.

\begin{lemma}[Exact selected ERM]
\label{lem:tie}
Let $S$ be the positive-weight support of a selected objective and let
$W_S=\spanop\{x_i:i\in S\}$. Its minimum-norm ERM equals $w_*$ if and only
if
\begin{equation}
 \sum_{i\in S}\lambda_i g_i=0
 \quad\text{and}\quad w_*\in W_S.
 \label{eq:tie-criterion}
\end{equation}
If $W_S=\R^d$, the second condition is automatic and $w_*$ is the unique
selected ERM.
\end{lemma}

The first condition says that $w_*$ minimizes the selected quadratic. Its
minimizer set is then $w_*+W_S^\perp$. The closest member to zero is $w_*$
exactly when $w_*\perp W_S^\perp$, proving the lemma. This distinction is
important. Translating $w_*$ to zero preserves loss differences but not the
minimum-norm choice on rank-deficient supports.

\paragraph{Positive bases.}
A finite set positively spans a vector space if every vector is a
nonnegative combination of its elements. An inclusion-minimal positive
spanning set is a positive basis. A positive basis of $\R^q$ has between
$q+1$ and $2q$ vectors. At the upper endpoint it is a cross:
after rescaling, it has the form $\{\pm u_1,\ldots,\pm u_q\}$ for a linear
basis $u_1,\ldots,u_q$ \citep{davis1954positive,regis2016positive}.

\section{One Example Below Exactness}

We first isolate the only configuration in which the standard $2d$ exact
certificate genuinely needs all $2d$ examples.

\begin{definition}[Unique-certificate number]
For a full-rank dataset, let $c_{\rm u}(D)$ be the least support size of a
weighted objective for which $w_*$ is the unique minimizer.
\end{definition}

\begin{theorem}[Cross dichotomy]
\label{thm:cross}
Every full-rank dataset satisfies $c_{\rm u}(D)\leq2d$. Equality holds if
and only if there are independent lines
$U_1,\ldots,U_d$ such that
\begin{enumerate}
\item every nonzero $g_i$ lies in some $U_j$,
\item both orientations occur among the $g_i$ on every $U_j$, and
\item $g_i=0$ implies $x_i=0$.
\end{enumerate}
If these conditions fail, $c_{\rm u}(D)\leq2d-1$.
\end{theorem}

\begin{proof}
Let $V=\spanop\{g_i:g_i\neq0\}$ and $q=\dim V$. The nonzero $g_i$ have a
strictly positive combination equal to zero. Steinitz's theorem therefore
selects at most $2q$ of them that positively span $V$. Every nonzero-gradient
feature belongs to $V$. Since all features span $\R^d$, we may add $d-q$
zero-gradient features that extend $V$ to $\R^d$. Their union gives a
unique exact objective on at most $d+q$ examples. Thus $q<d$ gives at most
$2d-1$.

Now let $q=d$ and choose a positive basis $B$ from the nonzero gradients.
If $|B|\leq2d-1$, its positive dependence and full span give the claim.
Otherwise $B$ is a cross $\{\pm u_j\}_{j=1}^d$. Write any other nonzero
gradient as
\begin{equation}
 g=\sum_{j\in J}a_j s_j u_j,
 \qquad a_j>0,\quad s_j\in\{-1,1\}.
\end{equation}
If $|J|\geq2$, then
\begin{equation}
 \{g\}\cup\{-s_j u_j:j\in J\}
 \cup\{\pm u_\ell:\ell\notin J\}
 \label{eq:replacement-cross}
\end{equation}
positively spans $\R^d$ and has $2d-|J|+1\leq2d-1$ elements. Hence every
nonzero gradient must lie on a cross line if $c_{\rm u}(D)=2d$.

Let $Z$ be the span of zero-gradient features and put $z=\dim Z$. If
$z>0$, choose $z$ such features and $d-z$ cross lines that together span
$\R^d$. A balancing pair from every chosen line plus the zero-gradient
features gives a unique exact objective on $2(d-z)+z=2d-z$ examples. Thus
$z=0$ is necessary.

Conversely, under the three stated conditions, full selected feature rank
requires using every line. Gradient balance holds independently on each
line and needs both orientations. Therefore at least $2d$ positive-weight
examples are necessary, while one balancing pair per line is sufficient.
\end{proof}

The theorem concerns uniqueness. By Lemma~\ref{lem:tie}, a rank-deficient
objective can still return $w_*$ when its feature span contains $w_*$. On
the exceptional branch, this tie behavior also has an exact description.

\begin{corollary}[Min-norm compression on a cross]
\label{cor:cross-compression}
Suppose the three conditions of Theorem~\ref{thm:cross} hold, and write
$w_*=\sum_{j=1}^d\beta_ju_j$. The least support size of a weighted
objective whose minimum-norm ERM is $w_*$ equals one if $w_*=0$ and the
dataset contains a zero-feature example. In every other case it equals
\begin{equation}
 2\max\{1,|\{j:\beta_j\neq0\}|\}.
 \label{eq:cross-compression}
\end{equation}
\end{corollary}

For every nonzero coordinate $\beta_j$, the selected feature span must
include $U_j$. Gradient balance on an independent cross line requires both
orientations. Conversely, balancing pairs on precisely those lines have a
minimizer set whose feature span contains $w_*$. If $w_*=0$, any one
balancing pair suffices, while a zero-feature loss has minimum-norm ERM zero
by itself. Thus $2d$ is necessary for the raw min-norm rule exactly when
every $\beta_j$ is nonzero.

\begin{theorem}[Sharp endpoint]
\label{thm:endpoint}
For every integer $d\geq1$,
\begin{equation}
 \sup_{D\subset\R^d\times\R}
 \frac{L_D^*(2d-1;\cA_*)}{L_D^*}=1+\frac1d.
 \label{eq:endpoint}
\end{equation}
\end{theorem}

\begin{proof}[Upper bound]
If the feature rank is $r<d$, the known $2r$ exact guarantee fits within
$2d-1$. The same is immediate when $L_D^*=0$. Assume full rank and positive
optimal loss. If Theorem~\ref{thm:cross} gives a unique exact support of
size at most $2d-1$, there is nothing to prove. We may therefore work under
its three cross conditions.

Choose a nonzero vector $u_j$ on each cross line and write
$x_i=\alpha_i u_j$ for indices $i\in I_j$. Put
\begin{equation}
 h_j=\frac1N\sum_{i\in I_j}\alpha_i^2,\qquad
 A_j=\frac1N\sum_{i\in I_j}r_i^2.
 \label{eq:line-shares}
\end{equation}
Examples with $x_i=0$ may add a constant to $L_D^*$ but do not enter these
sums. Let $q_1,\ldots,q_d$ be the dual basis, so
$u_\ell^\top q_j=\mathbf1\{\ell=j\}$. Equation~\eqref{eq:pythagorean}
gives
\begin{equation}
 (tq_j)^\top H(tq_j)=h_jt^2.
 \label{eq:dual-excess}
\end{equation}
Moreover $A_j/h_j$ is the $\alpha_i^2$-weighted average of
$(r_i/\alpha_i)^2$. Hence there is an $i(j)\in I_j$ with
\begin{equation}
 h_j\left(\frac{r_{i(j)}}{\alpha_{i(j)}}\right)^2\leq A_j.
 \label{eq:line-average}
\end{equation}
The line normal equation $\sum_{i\in I_j}r_i\alpha_i=0$ provides an example
with opposite gradient orientation to $i(j)$.

Pick a line $j$ with $A_j\leq L_D^*/d$. On every other line choose a
balancing pair. On line $j$ keep only $i(j)$. Weight each pair so its
gradient at $w_*$ cancels. The selected features span $\R^d$, and the
selected objective separates in the coordinates $u_\ell^\top(w-w_*)$.
Its unique minimizer is
\begin{equation}
 w_*+\frac{r_{i(j)}}{\alpha_{i(j)}}q_j.
\end{equation}
Equations~\eqref{eq:dual-excess} and \eqref{eq:line-average} bound its
full loss by $L_D^*+A_j\leq(1+1/d)L_D^*$.
\end{proof}

\begin{proof}[Lower bound]
Use the $2d$ examples
\begin{equation}
 (e_j,0),\quad(-e_j,-2),\qquad j=1,\ldots,d.
 \label{eq:cross-lower}
\end{equation}
The full ERM is $w_*=(1,\ldots,1)$ and $L_D^*=1$. Any support of size at
most $2d-1$ leaves some coordinate with at most one selected example. A
singleton makes that coordinate zero or two. With no selected example, the
minimum-norm rule makes it zero. Its two full-data losses therefore sum to
four instead of the optimal two. All other coordinates contribute at least
their optimal value, so the ratio is at least $1+1/d$. Equal weights on
complete pairs and one singleton attain equality.
\end{proof}

\begin{proposition}[Instance-dependent cross bound]
\label{prop:instance}
Under the cross conditions, define
\begin{equation}
 B_j=h_j\min_{i\in I_j}\left(\frac{r_i}{\alpha_i}\right)^2.
 \label{eq:instance-bound}
\end{equation}
There is a full-rank support of size $2d-1$ whose full-loss excess is
exactly $\min_jB_j$. Moreover,
\begin{equation}
 \min_jB_j\leq\min_jA_j\leq\frac{L_D^*}{d}.
\end{equation}
\end{proposition}

The construction in the upper-bound proof gives the first statement after
choosing the best line and singleton. The averaging argument
\eqref{eq:line-average} gives the inequalities. Thus a line containing an
example with a small residual-to-feature ratio can make the endpoint loss
much smaller than the universal $1+1/d$ bound.

\paragraph{Generic datasets are easier.}
If the nonzero gradients span $\R^d$ and are in linear general position,
an inclusion-minimal positive dependence has exactly $d+1$ elements and
spans $\R^d$. Thus $d+1$ examples recover $w_*$ uniquely. The intermediate
worst case is entirely a degeneracy phenomenon, with the endpoint extremum
at the maximally degenerate cross.

\paragraph{Constructive extraction.}
All supports in the proof can be found in polynomial time. Compute $w_*$,
the gradients $g_i$, and bases for their span and the zero-gradient feature
span. Repeatedly delete a gradient whenever the remainder still positively
spans the same space. Each test is a linear program, and the terminal set is
a positive basis of size at most twice its rank. A size-$2d$ terminal basis
is a cross. Equation~\eqref{eq:replacement-cross} either replaces it by a
smaller exact certificate, or \eqref{eq:instance-bound} chooses the
singleton. Positive balancing weights follow from one final linear
feasibility problem. The endpoint guarantee is therefore constructive, not
only existential.

\section{All-Budget Harmonic Lower Bound}

Write the intermediate budget as $n=d+k$ with $0\leq k<d$. For
$m\in\{k+1,\ldots,d\}$ define
\begin{equation}
 H(d,m)=\min_{\substack{b_j\in\mathbb N\\\sum_{j=1}^m b_j=d}}
 \sum_{j=1}^m\frac1{b_j}.
 \label{eq:harmonic}
\end{equation}
If $d=qm+r$ with $0\leq r<m$, convexity of $1/x$ gives
\begin{equation}
 H(d,m)=\frac{m-r}{q}+\frac{r}{q+1}.
 \label{eq:harmonic-closed}
\end{equation}
Set
\begin{equation}
 \Psi(d,k)=\max_{m=k+1,\ldots,d}\frac{m-k}{H(d,m)}.
 \label{eq:psi}
\end{equation}

\begin{theorem}[Intermediate lower bound]
\label{thm:lower}
For every $0\leq k<d$,
\begin{equation}
 \sup_D\frac{L_D^*(d+k;\cA_*)}{L_D^*}
 \geq 1+\Psi(d,k).
 \label{eq:general-lower}
\end{equation}
At $k=0$ this is $d+1$, and at $k=d-1$ it is $1+1/d$.
\end{theorem}

\begin{table}[t]
\caption{The lower-bound excess $\Psi(d,k)$ for small dimensions. The
selected budget is $d+k$.}
\label{tab:psi}
\centering
\small
\begin{tabular}{c@{\quad}cccccc}
\toprule
$d\backslash k$ & 0 & 1 & 2 & 3 & 4 & 5\\
\midrule
2 & $2$ & $1/2$ \\
3 & $3$ & $2/3$ & $1/3$ \\
4 & $4$ & $1$ & $1/2$ & $1/4$ \\
5 & $5$ & $6/5$ & $3/5$ & $2/5$ & $1/5$ \\
6 & $6$ & $3/2$ & $2/3$ & $1/2$ & $1/3$ & $1/6$ \\
\bottomrule
\end{tabular}
\end{table}

The construction uses orthogonal simplex blocks. For a composition
$d=b_1+\cdots+b_m$, block $j$ contains the standard $b_j+1$-point
simplex lower instance for size-$b_j$ regression. Give it optimal-loss
share
\begin{equation}
 a_j=\frac{1/b_j}{\sum_{\ell=1}^m1/b_\ell}.
 \label{eq:equalize}
\end{equation}
Using $b_j$ examples in that block incurs excess $b_ja_j$, while using all
$b_j+1$ removes the excess. A large parameter offset makes fewer than
$b_j$ examples no better under the minimum-norm rule. Every finite-ratio
selection spends the baseline $\sum_jb_j=d$ examples, and each of the $k$
remaining examples can complete one block. Equation~\eqref{eq:equalize}
equalizes all block excesses, leaving $(m-k)/\sum_j1/b_j$. The appendix
gives the finite-offset construction and limiting argument.

The lower bound has two regimes up to integer effects. The maximizing
$m$ is near $2k$ while $k\lesssim d/2$, giving excess about $d/(4k)$.
Once $m=d$ binds, the excess is $(d-k)/d$. The latter includes the endpoint
proved universally above.

\section{When the Harmonic Curve Is Exact}

The lower construction belongs to a broader class on which
\eqref{eq:general-lower} is also an upper bound. This result explains why
simplex blocks arise and states precisely which reduction is missing in the
general case.

Whiten the full loss around $w_*$ so that
\begin{equation}
 \sum_i p_ix_ix_i^\top=I,\quad
 \sum_i p_ir_ix_i=0,\quad
 \sum_i p_ir_i^2=1.
 \label{eq:normal-form}
\end{equation}
Then the full loss ratio at $w_*+v$ is $1+\|v\|^2$. Splitting an example
weight into nonnegative pieces does not change the objective and lets one
treat repeated circuit occurrences separately.

\begin{definition}[Rank-additive circuit decomposition]
\label{def:circuits}
A normalized instance is rank-additive circuit decomposable if its weighted
loss can be split into $m$ centered circuit quadratics
\begin{equation}
 F_j(v)=\|U_jv-s_j\|^2=a_j+v^\top H_jv,
 \quad H_j=U_j^\top U_j,
 \label{eq:circuit-quadratic}
\end{equation}
where $U_j$ has $b_j+1$ rows and rank $b_j$, $U_j^\top s_j=0$,
$a_j=\|s_j\|^2$, and
\begin{equation}
 \sum_jH_j=I,\quad \sum_ja_j=1,\quad \sum_jb_j=d.
 \label{eq:rank-additive}
\end{equation}
Every row occurrence must be a nonnegative scaling of an original example
loss.
\end{definition}

The condition arises from a conformal circuit decomposition of the full
residual in the kernel of the design. Rank additivity is the substantive
part. General decompositions can have overlapping circuits whose ranks sum
to more than $d$.

\begin{theorem}[Circuit sparsification]
\label{thm:circuits}
Under Definition~\ref{def:circuits}, for every $K\subseteq[m]$ there is a
weighted support of at most $d+|K|$ original examples with full-loss excess
at most
\begin{equation}
 \sum_{j\notin K}b_ja_j.
 \label{eq:circuit-bound}
\end{equation}
Consequently the worst ratio over rank-additive circuit decomposable
instances at budget $d+k$ is exactly $1+\Psi(d,k)$.
\end{theorem}

\begin{proof}
Choose full-column factors $H_j=B_jB_j^\top$, with
$B_j\in\R^{d\times b_j}$, and concatenate them as
$B=[B_1\ \cdots\ B_m]$. By \eqref{eq:rank-additive}, $B$ is square and
$BB^\top=I$. Thus $B$ is orthogonal. It follows that the columns within
each $B_j$ are orthonormal, columns from different factors are orthogonal,
and $H_j=B_jB_j^\top$ is the orthogonal projector onto its range.

Write $\widetilde U_j=U_jB_j$. Since $U_j$ vanishes on the orthogonal
complement of $\operatorname{range}(H_j)$, this gives the intrinsic
factorization $U_j=\widetilde U_jB_j^\top$. Moreover,
$\widetilde U_j^\top\widetilde U_j=I_{b_j}$,
$\widetilde U_j^\top s_j=0$, and
$\|s_j\|^2=a_j$. Hence the square augmented matrix
\begin{equation}
 [\,\widetilde U_j\ \ s_j/\sqrt{a_j}\,]
\end{equation}
is orthogonal when $a_j>0$.

Omit one of the $b_j+1$ circuit rows with probability proportional to the
squared volume of the remaining $b_j$ design rows. The resulting
interpolator $v_j$ has
\begin{equation}
 \E v_j=0,\qquad \E[v_jv_j^\top]=a_jH_j,\qquad
 \E\|v_j\|^2=b_ja_j.
 \label{eq:circuit-covariance}
\end{equation}
These identities follow directly by viewing the omitted row of the
orthogonal augmented matrix as its cofactor vector. They also follow from
the size-$b_j$ volume-sampling identity.

For $j\in K$, retain all $b_j+1$ rows and use its centered circuit weights.
For $j\notin K$, retain the random $b_j$-row basis. Orthogonality makes the
global selected ERM the sum of the incomplete-circuit interpolators. Thus
its expected squared norm is the right side of
\eqref{eq:circuit-bound}. Some deterministic choice is no worse. It uses
$\sum_jb_j+|K|=d+|K|$ row occurrences and no more original examples.

If $m\leq k$, the selector completes every circuit and is exact. Otherwise,
for fixed $b_j$, it completes the $k$ largest values $b_ja_j$. The exact
allocation identity for $m\geq k+1$ is
\begin{equation}
 \sup_{a\geq0,\ \sum a_j\leq1}
 \min_{|K|=k}\sum_{j\notin K}b_ja_j
 =\max_{\substack{J\subseteq[m]\\|J|\geq k+1}}
 \frac{|J|-k}{\sum_{j\in J}1/b_j}.
 \label{eq:minimax}
\end{equation}
It is proved in the appendix. The active subset $J$ allows zero loss on
some circuit blocks. Optimizing also over the number and integer ranks of
the circuits reduces the right side to $\Psi(d,k)$. Theorem~\ref{thm:lower}
supplies equality.
\end{proof}

\section{Discussion}

The endpoint $2d-1$ is governed by a sharp geometric dichotomy. A single
gradient involving two cross coordinates creates a smaller positive basis.
A single nonzero zero-gradient feature also lowers the exact unique-ERM
certificate. Only the cross survives, and its loss allocation forces the
$1/d$ constant. This is why the endpoint admits an exact argument despite
the mismatch between volume sampling and Steinitz sparsification.

The all-budget lower bound suggests the curve $1+\Psi(d,k)$. Theorem
\ref{thm:circuits} proves it whenever conformal circuits are rank additive.
Positive-basis decompositions need not have this property. Their simplices
may overlap through critical vectors \citep{schoch2020simplicial,
hare2021positive}. Such overlap gives additional cancellation routes in
small-dimensional exact optimization, but turning that observation into a
universal inequality remains the central open step.

Regularized volume sampling controls inverse Hessians under stochastic
noise assumptions \citep{derezinski2018ridge}. It does not control the
fixed residual-induced bias here. Ordinary size-$s>d$ volume sampling also
lacks a response-uniform loss identity \citep{derezinski2018reverse}. The
circuit proof succeeds because each local residual space is one-dimensional
and rank additivity makes its covariance orthogonal to every other circuit.

\bibliography{references}

\appendix

\section{Additional Linear-Algebra Details}

\begin{proof}[Expanded proof of Lemma~\ref{lem:tie}]
The selected objective has gradient at $w_*$ equal to
$-2\sum_{i\in S}\lambda_i g_i$. Therefore its first condition for
optimality is exactly the first part of \eqref{eq:tie-criterion}. Under this
condition its value at $w_*+v$ differs from its value at $w_*$ by
\begin{equation}
 \sum_{i\in S}\lambda_i(x_i^\top v)^2.
\end{equation}
Its minimizer set is $w_*+W_S^\perp$. The minimum-norm member of this affine
space is the orthogonal projection of $w_*$ onto $W_S$. It equals $w_*$ if
and only if $w_*\in W_S$.
\end{proof}

\begin{lemma}[Cross replacement]
\label{lem:replacement}
Let $B=\{\pm u_1,\ldots,\pm u_d\}$ be a cross. If a nonzero vector $g$
has at least two nonzero coordinates in the basis $(u_j)$, then
$B\cup\{g\}$ contains a positive spanning subset of size at most $2d-1$.
\end{lemma}

\begin{proof}
Write $g=\sum_{j\in J}a_js_ju_j$ with $a_j>0$ and $s_j\in\{-1,1\}$.
The set $\{g\}\cup\{-s_ju_j:j\in J\}$ positively spans the coordinate
subspace indexed by $J$, because
$g+\sum_{j\in J}a_j(-s_ju_j)=0$ is a strictly positive dependence and the
set spans that subspace. Add both orientations on every coordinate outside
$J$. The union positively spans $\R^d$ and has
$1+|J|+2(d-|J|)=2d-|J|+1\leq2d-1$ elements.
\end{proof}

\section{Finite Lower-Bound Construction}
\label{app:lower}

Fix a block dimension $b$ and an offset $M>0$. Use features
\begin{equation}
 e_1,\ldots,e_b,-\mathbf1_b
\end{equation}
and responses
\begin{equation}
 M,2M,\ldots,bM,
 -M\frac{b(b+1)}2-b-1.
 \label{eq:block-responses}
\end{equation}
The full minimizer is
\begin{equation}
 w_*=(M+1,2M+1,\ldots,bM+1),
\end{equation}
and every one of the $b+1$ residuals has magnitude one. The average optimal
loss is one. Omitting any one point and interpolating the other $b$ makes
the omitted residual have magnitude $b+1$. Hence the full average loss is
$b+1$, with excess $b$.

Consider any support with fewer than $b$ distinct examples. Its selected
rows are linearly independent, so every positive weighting has zero selected
loss and returns the minimum-norm interpolator of those rows. If
$-\mathbf1_b$ is absent, at least one omitted basis coordinate is set to
zero, while its response has magnitude at least $M$. The full loss therefore
diverges with $M$.

If $-\mathbf1_b$ is present, at most $b-2$ basis rows are present. Let $J$
be the set of at least two omitted basis coordinates. The minimum-norm
solution assigns a common value to every coordinate in $J$, because their
only constraint is their sum. Two distinct targets in $J$ differ by at
least $M$. For any common prediction $c$ and two such indices $q\ne r$,
\begin{equation}
 (c-qM)^2+(c-rM)^2\geq \frac{(q-r)^2M^2}{2}\geq\frac{M^2}{2}.
\end{equation}
Thus the full loss again diverges as $M\to\infty$. This is the same offset
device used for the sharp size-$d$ regression lower bound of
\citet{hanneke2025data}.

Take orthogonal copies with dimensions $b_1,\ldots,b_m$. Scale all responses
in block $j$ so its contribution to the full optimal loss is $a_j$. Orthogonal
features make both the selected ERM and full loss decompose blockwise. In
the limit $M\to\infty$, a finite-ratio support uses at least $b_j$ examples
in every block. With exactly $b_j$ its excess contribution is $b_ja_j$.
With $b_j+1$ it is zero. A budget $d+k$ completes at most $k$ blocks.
Choosing \eqref{eq:equalize} proves Theorem~\ref{thm:lower}. For every finite
$M$ this is a finite dataset, so taking $M\to\infty$ is valid for the
supremum in \eqref{eq:general-lower}.

\section{The Allocation Identity}

\begin{proof}[Proof of \eqref{eq:minimax}]
Put $c_j=b_ja_j$ and let $T(c)$ be the sum of the $m-k$ smallest $c_j$.
For any active set $J$ with $|J|\geq k+1$, set $c_j=t$ on $J$ and zero
outside, where $t=(\sum_{j\in J}1/b_j)^{-1}$. This is feasible and gives
$T(c)=(|J|-k)t$, proving the lower bound in \eqref{eq:minimax}.

For the reverse bound use the layer-cake identities
\begin{align}
 T(c)&=\int_0^\infty (|J_t|-k)_+\,dt,\\
 \sum_j\frac{c_j}{b_j}
 &=\int_0^\infty\sum_{j\in J_t}\frac1{b_j}\,dt,
 \end{align}
where $J_t=\{j:c_j\geq t\}$. If the right side of
\eqref{eq:minimax} is $\Phi$, its definition gives
$(|J_t|-k)_+\leq\Phi\sum_{j\in J_t}1/b_j$ at every $t$. Integration gives
$T(c)\leq\Phi\sum_jc_j/b_j\leq\Phi$.
\end{proof}

To optimize over compositions, fix an active subset $J$ of size $\ell$.
Since $\sum_{j\in J}b_j\leq d$, adding unused dimensions to these blocks can
only decrease $\sum_{j\in J}1/b_j$. Hence
\begin{equation}
 \sum_{j\in J}\frac1{b_j}\geq H(d,\ell).
\end{equation}
The corresponding allocation value is at most
$(\ell-k)/H(d,\ell)\leq\Psi(d,k)$. Conversely, take $m=\ell$, use the
balanced composition attaining $H(d,\ell)$, and activate every block. This
proves that the worst fixed-composition value, optimized over all
compositions, is exactly $\Psi(d,k)$.

\section{Circuit Covariance Identity}

Let $U\in\R^{(b+1)\times b}$ satisfy $U^\top U=I_b$, and let
$s\in\ker(U^\top)$ have $\|s\|^2=a>0$. Then
\begin{equation}
 Q=[\,U\ \ s/\sqrt a\,]
\end{equation}
is an orthogonal matrix. For omission $i$, let $U_{-i}$ and $s_{-i}$ remove
row $i$, and put $z=s/\sqrt a$. Cofactor identities give the determinant
formula below. Orthogonality of row $i$ with every remaining row gives
$U_{-i}u_i+z_{-i}z_i=0$, and therefore
\begin{equation}
 \det(U_{-i})^2=\frac{s_i^2}{a},\qquad
 U_{-i}^{-1}s_{-i}=-\sqrt a\,\frac{u_i}{s_i/\sqrt a},
 \label{eq:cofactor}
\end{equation}
where $u_i^\top$ is row $i$ of $U$. Consequently volume sampling has
probabilities $s_i^2/a$ and
\begin{align}
 \E[U_{-i}^{-1}s_{-i}]&=0,\\
 \E[(U_{-i}^{-1}s_{-i})(U_{-i}^{-1}s_{-i})^\top]
 &=a\sum_i u_iu_i^\top=aI_b.
\end{align}
The mean identity is also the standard unbiasedness of size-$b$ volume
sampling. Mapping intrinsic coordinates back to the range of $H_j$ proves
\eqref{eq:circuit-covariance}. The case $a=0$ follows by continuity and has
zero excess.

\end{document}
